% Captions of the figures in figures/, made by make_figures.py. Every number quoted here is in figure_metrics.json;
% if a figure changes, update both. Notation as in closed_form.tex.
% "Validation proteins": the 695 validation proteins of the cluster split (MMseqs2 clusters at 30% identity and 80%
% coverage; no validation protein shares a cluster with a training protein). MD: HPS1 simulations throughout.

\begin{figure}
\centering
\includegraphics[width=\columnwidth]{fig0_intro_example}
\caption{\textbf{The network predicts the full distribution of the radius of gyration.}
Predicted $P(R_g)$ of a validation protein ($N=95$); dashed line: its mean $\langle R_g\rangle=2.56$\,nm (MD: 2.53\,nm).}
\label{fig:intro}
\end{figure}

% fig0b replaces fig0 in the paper
\begin{figure}
\centering
\includegraphics[width=\columnwidth]{fig0b_intro_pair}
\caption{\textbf{The network predicts the full distribution of the radius of gyration.}
Predicted $P(R_g)$ of two validation proteins with the same mean, $\langle R_g\rangle=4.00$ and 4.03\,nm (MD: 3.97 and
3.99\,nm; dotted line), but different exponents $\nu$: the longer chain ($N=189$, $\nu=0.59$) has a broader,
more skewed distribution than the shorter one ($N=96$, $\nu=0.73$).}
\label{fig:intro_pair}
\end{figure}

\begin{figure}
\centering
\includegraphics[width=\columnwidth]{fig1_val_moments}
\caption{\textbf{Moments of the predicted distributions against MD.}
Mean $\langle R_g\rangle$ (\textbf{a}), standard deviation $\sigma_{R_g}$ (\textbf{b}), skewness (\textbf{c}) and excess
kurtosis (\textbf{d}) of the network's $P(R_g)$ against those of the MD frames, for the 695 validation proteins. Dashed
line: equality. $R^2$ and mean absolute error (MAE) are computed about the equality line.}
\label{fig:val_moments}
\end{figure}

\begin{figure}
\centering
\includegraphics[width=\columnwidth]{fig24_val_nu}
\caption{\textbf{The exponent against MD.}
The network's exponent $\nu$ against the one fitted by maximum likelihood to each protein's MD frames, for the 695
validation proteins. Dashed line: equality. $R^2$ and mean absolute error (MAE) are computed about the equality line.}
\label{fig:val_nu}
\end{figure}

\begin{figure}
\centering
\includegraphics[width=\columnwidth]{fig2_val_distributions}
\caption{\textbf{Predicted distributions against MD, from typical to worst.}
\textbf{a},~Wasserstein-1 distance $W_1$ between the network's $P(R_g)$ and the distribution of the MD frames for the 695
validation proteins, as a percentage of the MD $\langle R_g\rangle$ ($W_1$ is the area between the two cumulative
distributions); median 1.5\%. Red lines: the proteins in \textbf{b}--\textbf{e}.
\textbf{b}--\textbf{e},~MD frames (histogram) and the network's $P(R_g)$ of the proteins at the 25th (0.9\%), 50th
(1.5\%) and 75th (2.5\%) percentiles of $W_1$, and of the protein with the largest $W_1$ (14.6\%, a compact chain).}
\label{fig:val_distributions}
\end{figure}

\begin{figure*}
\centering
\includegraphics[width=\textwidth]{fig4_real_proteins}
\caption{\textbf{Predicted distributions of the benchmark proteins.}
The benchmark: 122 proteins with a SAXS $R_g$ and an HPS1 trajectory, none used in training or model selection.
\textbf{a},~$W_1$ between the network's $P(R_g)$ and the distribution of each protein's HPS1 trajectory, as a percentage of
its MD $\langle R_g\rangle$ (median 1.6\%, 90th percentile 3.4\%, largest 7.9\%). Solid lines mark the proteins shown in
\textbf{b--e}; dashed line: the median.
\textbf{b--e},~Proteins at the 10th, 50th and 90th percentile and the largest $W_1$: histogram of the MD frames and the
network's prediction.}
\label{fig:benchmark}
\end{figure*}

\begin{figure}
\centering
\includegraphics[width=\columnwidth]{fig13_experiment_parity}
\caption{\textbf{Mean radius of gyration against experiment.}
$\langle R_g\rangle$ of the 122 benchmark proteins from HPS1 simulations (circles) and from the network (triangles), against
SAXS. Colours: the three largest sequence families (proteins linked by sequence similarity $\geq0.6$ or the same
amino-acid composition); all other proteins in green. The network reproduces the simulations it was trained on (MAE 0.059\,nm against
HPS1) and so inherits their deviation from experiment: MAE against SAXS 0.42\,nm for the network and 0.39\,nm for HPS1.
Dashed line: equality.}
\label{fig:saxs}
\end{figure}

\begin{figure*}
\centering
\includegraphics[width=\textwidth]{fig15_surrogate_examples}
\caption{\textbf{Closed form against the network, typical to worst.}
Validation proteins at the median (\textbf{a}) and 95th percentile (\textbf{b}) of the Kullback--Leibler divergence
$\mathrm{KL}(p_\mathrm{network}\,\|\,p_\mathrm{closed\,form})$ over the 695 validation proteins, and the protein with the
largest divergence (\textbf{c}): MD frames (histogram), network (solid) and closed form (dashed). The largest divergence is
a compact chain that both miss: MD $\langle R_g\rangle=3.05$\,nm, network 2.68\,nm, closed form 2.39\,nm.}
\label{fig:cf_examples}
\end{figure*}

\begin{figure*}
\centering
\includegraphics[width=\textwidth]{fig20_surrogate_summary}
\caption{\textbf{How closely the closed form reproduces the network.}
\textbf{a--c},~As Fig.~\ref{fig:cf_examples}.
\textbf{d},~$\mathrm{KL}(p_\mathrm{network}\,\|\,p_\mathrm{closed\,form})$ over the 695 validation proteins: median 0.0019
nats, 95th percentile 0.031, 99th percentile 0.12; the last bin collects the 31 proteins with $\mathrm{KL}\geq0.04$.
\textbf{e},~The exponent $\nu$ of the two-exponent distribution: fitted by maximum likelihood to each protein's MD frames
(filled), predicted by the network (solid) and by the closed form (dashed). Against the fitted values, $R^2=0.72$ for the
network and 0.67 for the closed form. The closed form's $\nu$ spans 0.41--0.89, the same range as the fitted values and wider
than the network's (0.41--0.82).
Dashed vertical lines: $\theta$ chain ($\nu=1/2$) and good solvent ($\nu=0.588$).}
\label{fig:cf_summary}
\end{figure*}

\begin{figure*}
\centering
\includegraphics[width=\textwidth]{fig21_surrogate_response}
\caption{\textbf{Response of the closed form to its inputs.}
Each row varies one input of the closed form over the 1st--99th percentile range of the 3{,}223 training proteins,
with the other inputs held at a reference chain (black circle): $N=184$ and $\langle\lambda\rangle=0.603$ (training
medians), contact energy at the CE threshold $c(N)=c_0+0.25\ln N$ ($-0.128$ at $N=184$; along \textbf{a} the contact
energy follows $c(N)$), so that $h=0$, and no charge decoration ($s_+=s_-=0$); its $\langle R_g\rangle=4.18$\,nm,
$\nu=0.606$ and $A=3.33$. Left: change of the mean $\langle R_g\rangle$ (solid), the exponent $\nu$ (dashed) and the
amplitude $A$ (dotted) of the predicted distribution relative to the reference chain; \textbf{b--e} share their vertical
axis. Right: $P(R_g)$ at $N=184$ of the chain marked by the diamonds, with one input at the end of its training range
(1st percentile for $s_-$, 99th for the others), and of the reference chain (grey outline), all on one density scale;
short vertical line: the chain's mean; dashed line: the reference mean.
\textbf{a},~Chain length, on a logarithmic axis: $\langle R_g\rangle$ rises from 1.27\,nm at $N=24$ to 6.68\,nm at
$N=427$, while $\nu$ stays at 0.606, since at $s_+=0$ the exponent does not depend on $N$, and $A$ falls with the $1/N$ term of
its law (5.05 to 3.21).
\textbf{b},~$s_-=\min[\operatorname{asinh}(5\,\mathrm{SCD}/N),0]$, negative when opposite charges are segregated into
blocks; at $s_-=-0.255$, $\langle R_g\rangle$ $-11\%$, $\nu=0.540$, $A=2.84$.
\textbf{c},~$s_+=\max[\operatorname{asinh}(5\,\mathrm{SCD}/N),0]$, positive when like charges dominate; at $s_+=2.25$,
$\langle R_g\rangle$ $+48\%$, $\nu=0.702$, $A=4.17$.
\textbf{d},~Contact energy CE, which enters only above the CE threshold (dotted line, $c(184)$) through
$h=\max[\mathrm{CE}-c(N),0]$; at $\mathrm{CE}=2.75$, $\langle R_g\rangle$ $-18\%$, $\nu=0.507$, $A=2.63$.
\textbf{e},~Mean HPS hydropathy $\langle\lambda\rangle$; at
$\langle\lambda\rangle=0.766$, $\langle R_g\rangle$ $-6\%$, $\nu=0.603$, $A=3.13$.}
\label{fig:cf_response}
\end{figure*}

\begin{figure*}
\centering
\includegraphics[width=\textwidth]{fig23_shape_from_sequence}
\caption{\textbf{Proteins with the same mean radius of gyration can have different distributions, and the network
predicts which.}
Pairs of validation proteins whose MD $\langle R_g\rangle$ agree within 1\%, regardless of chain length (3{,}558 pairs).
In 1{,}177 of these pairs the MD skewness differs by at least 0.2; the network orders the two skewnesses correctly in
96\% of them (closed form: 96\%). Shown are the four of those pairs whose MD distributions differ most (Wasserstein-1
distance 6.1--6.5\% of $\langle R_g\rangle$), with no protein used twice; the pairs were selected from the MD data alone.
In each panel the less skewed protein is drawn upward (blue) and the other downward (orange, mirrored): histogram of its
MD frames and the density the network predicts from its sequence. Dashed line: the mean of the two MD
$\langle R_g\rangle$. Labels: chain length $N$, the network's exponent $\nu$, and the MD standard deviation $\sigma$ and
skewness, with the network's values in parentheses. In every pair a short, expanded chain ($\nu=0.71$--$0.82$) is
matched with a longer, more compact one ($\nu=0.44$--$0.60$).}
\label{fig:shape_pairs}
\end{figure*}

\begin{figure*}
\centering
\includegraphics[width=\textwidth]{fig25_family_heads}
\caption{\textbf{The two-exponent output against free mixture densities.}
The network retrained with only its output head changed: log-normal mixtures with $K$ components (blue circles) and
two-exponent forms (red squares), same inputs, NLL only, last epoch. Numbers in parentheses: free output parameters per
protein; the released head (bottom row, dotted line) has 2, with six global coefficients for $A$. Networks are
$512\times2$ unless stated; the open circle is the earlier recipe, with moment penalties and checkpoint selection on
validation. Mean and sd over three seeds (open circle on the sticky split: two seeds).
\textbf{a}, Validation NLL: the released head has 1.1524, the $K=3$ network at $512\times3$ 1.1508.
\textbf{b}, Mean Wasserstein-1 distance on the sticky split, sequences more contact-prone than any used in training
(5.96\% against 6.67\% for that network).
\textbf{c}, The same distance against network size (labels: width $\times$ hidden layers). At $64\times2$ the mixture
reaches 11.42\%, the two-exponent head 7.38\%.}
\label{fig:family_heads}
\end{figure*}

\begin{figure*}
\centering
\includegraphics[width=\textwidth]{fig26_starling_shape}
\caption{\textbf{Shape of the $R_g$ distribution: the network against STARLING.}
The network is trained on HPS1 and STARLING on Mpipi-GG. Every distribution is divided by its own mean, so size plays no part.
\textbf{a--c}, Spread $\sigma/\langle R_g\rangle$ and the probabilities of compact ($R_g<t_c\langle R_g\rangle$) and extended
($R_g>t_e\langle R_g\rangle$) conformations, where $t_c=0.681+0.311/\sqrt{N}$ and $t_e=1.421-0.571/\sqrt{N}$ are the 5th and 95th percentiles of
$R_g/\langle R_g\rangle$ of a typical training chain of the same length. From the network against STARLING (blue circles) and, for scale, against the HPS1
trajectories (green triangles), the force field the network was trained on. Grey band: $\pm1.96$ times the median sampling
error of a STARLING ensemble (2000 conformations); block resampling puts HPS1's slightly lower. $\rho$: Spearman correlation and MAE over the
119 proteins with $N\geq20$; open symbols have $N<20$, outside the network's domain.
\textbf{d--g}, The proteins at the 25th, 50th and 75th percentiles and the largest of the shape $W_1$ between the network and
STARLING over 121 proteins, as a percentage of $\langle R_g\rangle$: STARLING histogram and network prediction. Dotted lines mark $t_c$ and $t_e$ at
that length.}
\label{fig:starling_shape}
\end{figure*}
